% ===================================================================
%  L4: A Functional Programming Language for Computational Law
%  ICAIL submission DRAFT.
%
%  Build:  pdflatex l4-icail ; bibtex l4-icail ; pdflatex l4-icail ; pdflatex l4-icail
%  (or:    latexmk -pdf l4-icail.tex)
%
%  EDITORIAL TODOs (search "TODO"):
%   - co-authors and affiliations
%   - confirm what may be disclosed about each pilot (Sec. Experience)
%   - re-verify Hvitved/CSL operator correspondence against the thesis
% ===================================================================
\documentclass[sigconf,nonacm]{acmart}

% --- Packages -------------------------------------------------------
\usepackage{listings}
\usepackage{xcolor}
\usepackage{booktabs}

\settopmatter{printacmref=false}
\renewcommand\footnotetextcopyrightpermission[1]{}
\setcopyright{none}

% --- L4 listings style ----------------------------------------------
\definecolor{l4kw}{RGB}{0,68,136}      % keywords  (deep blue)
\definecolor{l4reg}{RGB}{150,30,120}   % regulative keywords (magenta)
\definecolor{l4cmt}{RGB}{110,110,110}  % comments
\definecolor{l4str}{RGB}{0,120,60}     % strings
\definecolor{l4bg}{RGB}{248,248,248}

\lstdefinelanguage{L4}{
  morekeywords={GIVEN,GIVES,GIVETH,MEANS,DECIDE,DECLARE,HAS,IS,ONE,OF,A,AN,THE,
    WITH,WHERE,CONSIDER,WHEN,THEN,ELSE,OTHERWISE,MATCHES,BRANCH,IF,AND,OR,NOT,
    TRUE,FALSE,IMPORT,ASSUME,LIST,EMPTY,FOLLOWED,BY,MAYBE,JUST,NOTHING,TYPE,
    PLUS,MINUS,TIMES,DIVIDED,MODULO,AT,LEAST,MOST,GREATER,THAN,LESS,EQUALS,APPEND},
  morekeywords=[2]{PARTY,MUST,MAY,SHANT,DO,WITHIN,BEFORE,AFTER,UPON,HENCE,LEST,
    BREACH,FULFILLED,BECAUSE,PROVIDED,EXACTLY,RAND,ROR,UPDATE,DOES},
  morekeywords=[3]{@export,@desc},
  sensitive=true,
  morecomment=[l]{--},
  morestring=[b]",
}

\lstset{
  language=L4,
  basicstyle=\ttfamily\footnotesize,
  keywordstyle=\color{l4kw}\bfseries,
  keywordstyle=[2]\color{l4reg}\bfseries,
  keywordstyle=[3]\color{l4kw}\itshape,
  commentstyle=\color{l4cmt}\itshape,
  stringstyle=\color{l4str},
  backgroundcolor=\color{l4bg},
  frame=tb,
  framesep=4pt,
  rulecolor=\color{l4cmt},
  columns=fullflexible,
  keepspaces=true,
  showstringspaces=false,
  upquote=true,
  breaklines=true,
  breakatwhitespace=true,
  literate={`}{{\char96}}1 {§}{{\S}}1 {'}{{\char39}}1 {>=}{{$\geq$}}1,
  aboveskip=6pt,
  belowskip=4pt,
}
% Inline code: texttt-based so it is robust inside macro arguments
% (lstinline cannot be safely captured as a macro argument).
\newcommand{\bq}{\char96\relax}            % straight backtick
\newcommand{\code}[1]{\texttt{\small #1}}

% --- Metadata -------------------------------------------------------
\begin{document}

\title{L4: A Functional Programming Language for Computational Law}
\subtitle{Constitutive Rules as Total Functions, Regulative Contracts as Concurrent Processes}

% TODO: co-authors and affiliations
\author{Wong Meng Weng}
\affiliation{%
  \institution{Centre for Computational Law, Singapore Management University; and Legalese}
  \city{Singapore}
  \country{}
}
\email{mengwong@legalese.com}

\author{The L4 Team}
\affiliation{%
  \institution{Centre for Computational Law, Singapore Management University}
  \city{Singapore}
  \country{}
}

\renewcommand{\shortauthors}{Wong et al.}

\begin{abstract}
We present \emph{L4}, a statically typed functional programming language
for computational law, together with its semantic foundations and the
controlled-natural-language surface syntax that makes it usable by legal
domain experts. L4 cleanly separates two readings of legal text that the
literature has long distinguished but that prior formalisms tend to
conflate. \emph{Constitutive} rules---what counts as what---are encoded as
total typed functions in the lineage of Haskell-style functional
programming, giving type safety, exhaustive case analysis, and an
audit-grade evaluation trace for every decision. \emph{Regulative}
rules---who must do what, by when, and what happens otherwise---are encoded
in a deontic/temporal layer whose denotation is a set of timed event
traces, borrowing heavily from Hvitved's Contract Specification Language
(CSL) and, through it, from Hoare's Communicating Sequential Processes and
the event calculus.
The same source program admits a second, relational reading via the
standard functional-to-relational transformation familiar from Mercury and
Curry; this lets L4 be evaluated forward as a decision procedure, run
backward for abductive planning, and exported as a transition relation to
symbolic model checkers, so that the search for legal loopholes becomes the
search for reachable ``double-bind'' states. We describe the language, its
implementation (parser, type checker, evaluator, language server, REST/MCP
decision service, and multiple compilation backends), and our design
rationale for a low-punctuation syntax---indentation in place of
parentheses, ditto marks, and phrase-shaped identifiers---intended to let
non-programmers read and write law in black and white.
\end{abstract}

\keywords{computational law, rules as code, low-code, controlled natural language,
functional programming, deontic logic, process calculi, contract
formalisation, legal knowledge representation}

\maketitle

% ===================================================================
\section{Introduction}
% ===================================================================
The corporation, viewed from a sufficient distance, is the sum of its
contracts: with shareholders, lenders, customers, vendors, employees, and
the state. Yet the legal department---custodian of this most fundamental
corporate asset---works in tooling that has barely advanced past the word
processor. Where the graphic designer has Adobe's suite and the programmer
commands languages, compilers, version control, debuggers, and test
frameworks, the lawyer has Microsoft Word. This is a striking asymmetry,
because the work is not so different: both the programmer and the drafter
translate a tangle of business goals into a precise textual artefact that
must anticipate the behaviour of a complex system with many moving parts
and adversarial participants.

This paper reports on \emph{L4}, a programming language built to close that
gap. L4 is not a document format, a clause library, or a contract-lifecycle
database; it is a language, in the sense that PostScript was a language
before desktop publishing was a product. Its design rests on three
commitments that, taken together, distinguish it from the existing body of
work in AI and Law. First, it takes seriously the classical distinction
between \emph{constitutive} and \emph{regulative} rules
\cite{searle1969speech,hart1961concept}, giving each a dedicated and
appropriate formal treatment rather than forcing both into a single
logic. Second, it grounds its regulative layer in a trace-based,
concurrency-aware semantics descended from contract DSLs and process
calculi \cite{hvitved2011thesis,hoare1978csp}, so that multi-party
obligations with deadlines and reparations are first-class. Third, it is
deliberately engineered as a \emph{controlled natural language}
\cite{kuhn2014cnl}: its concrete syntax minimises punctuation, replaces
parentheses with indentation, and admits identifiers that are whole legal
phrases, so that a statute and its formalisation can be read side by side
and recognised as the same text.

\paragraph{Contributions.} We make the following contributions.
\begin{itemize}
\item We describe the design of L4 and the rationale for its two-layer
  architecture: a total typed functional core for constitutive rules
  (Section~\ref{sec:constitutive}) and a deontic/temporal layer for
  regulative rules (Section~\ref{sec:regulative}).
\item We give the semantic foundations of the regulative layer as a
  trace-based denotational model, and exhibit the correspondence between
  L4's surface keywords (\code{HENCE}, \code{LEST}, \code{WITHIN},
  \code{RAND}, \code{ROR}) and the operators of Hvitved's CSL and Hoare's
  CSP (Section~\ref{sec:regulative}).
\item We show how a single L4 program supports two readings---functional
  and relational---via the function-as-graph transformation used in
  Mercury and Curry, and how the relational reading underwrites abductive
  planning queries and model-checking-based loophole detection
  (Section~\ref{sec:relational}).
\item We document the controlled-natural-language design of L4's surface
  syntax and the empirical motivation for each choice
  (Section~\ref{sec:cnl}).
\item We report on the implementation and on early deployments in the
  public and private sectors (Sections~\ref{sec:impl}
  and~\ref{sec:experience}).
\end{itemize}

% ===================================================================
\section{Motivation: Why a Language, in the Age of LLMs?}\label{sec:motivation}
% ===================================================================
A fair objection greets any new programming language in 2026: why bother?
Large language models can already program in every language that matters,
and a generation of ``vibe coders'' has, in effect, stopped writing code by
hand, describing intent in prose and letting the model fill in the
syntax. If the machine will write the code, why design a new notation for
humans to read?

The answer is that law is not like ordinary software, because of what
happens when it bites. Most of the time a citizen is content to remain a
rational ignoramus about the rules that govern them: the expected cost of
understanding a tax schedule or an insurance policy exceeds the expected
benefit, so people sign and move on. But the calculus inverts the moment
the stakes become concrete and personal---when a person might go to jail,
or opens an envelope expecting a bill for \$600 and finds one for
\$6{,}000. In that moment, ordinary people discover an enormous and
entirely rational capacity for detail. They want to understand the rule
\emph{themselves}; they are no longer willing to defer to a black box,
whether that box is a bureaucrat, a chatbot, or an opaque model whose
confident paraphrase may or may not track the actual entitlement. A
plausible-sounding generated answer is precisely what they cannot trust,
because the disagreement is now adversarial and the counterparty has every
incentive to read the words differently.

What such situations demand is a rule that is legible \emph{twice over}:
readable by a non-technical human who needs to check the reasoning against
their own circumstances, and readable by a computer that can evaluate it,
test it, and explain it without hallucinating. This dual-legibility
requirement is not new; it is the recurring ambition of what we would now call
\emph{low-code}. It is exactly the aspiration that gave us COBOL,
whose designers wanted business rules expressed in something a manager
could read, and SQL, whose declarative surface was meant to let analysts,
not just engineers, ask questions of their data. Both bet that some rules
are important enough that people will always insist on having them written
down in black and white. L4 inherits that bet and adds one refinement: the
black and white should be a \emph{formal} language. Natural-language
drafting carries avoidable, self-inflicted defects---lexical ambiguity,
scope ambiguity in coordinated clauses, the notorious unreliability of
punctuation\footnote{The missing-comma cases are a genre of their own;
\emph{O'Connor v.\ Oakhurst Dairy} turned a \$5M overtime dispute on the
absence of a serial comma.}---that a designed notation can simply
forbid. The role of the LLM, in our architecture, is not to replace this
notation but to help author and explain it: to draft a first formalisation
from source text, and to verbalise a rigorous evaluation trace back to the
user in accessible prose. The neural ``right brain'' handles language; a
deterministic ``left brain'' handles the rule. L4 is the left brain.

% ===================================================================
\section{Background and Design Goals}\label{sec:background}
% ===================================================================
\paragraph{Constitutive versus regulative.}
Searle distinguishes regulative rules, which govern antecedently existing
behaviour (``drive on the left''), from constitutive rules, which create
the very possibility of that behaviour (``this configuration of pieces
\emph{counts as} checkmate'') \cite{searle1969speech}. Hart's primary and
secondary rules and Hohfeld's jural relations draw related lines
\cite{hart1961concept,hohfeld1913}. The distinction is not academic for a
formaliser. A constitutive rule---``a person born in the United Kingdom
after commencement \emph{is} a British citizen if \ldots''---is a
definition: it relates facts to legal classifications and is naturally a
\emph{function} from situations to conclusions. A regulative rule---``the
buyer \emph{must} pay within thirty days, failing which \ldots''---is about
action over time, blame, and remedy, and is naturally a \emph{process}. The
two demand different mathematics. A recurring weakness of single-formalism
approaches is that one of the two is bent to fit the other: deontic notions
are encoded as ordinary predicates and lose their temporal and reparational
structure, or definitional computation is shoe-horned into a modal logic
where it does not belong. L4's central architectural decision is to give
each its own layer.

\paragraph{Isomorphism.}
Following the tradition of the British Nationality Act as a logic program
\cite{sergot1986bna} and the broader rules-as-code movement
\cite{mohun2020rulesascode}, L4 aims for \emph{isomorphic} encoding: the
structure of the formal text should mirror the structure of the source. If
a section has paragraphs (a), (b), and (c) joined by ``and'' and ``or'', the
encoding should exhibit the same conjunction-and-disjunction tree, so that a
domain expert can audit the correspondence clause by clause and so that a
later amendment to the source maps to a localised change in the
code. Isomorphism is simultaneously a correctness discipline and a
maintenance strategy; it also drives several surface-syntax choices
(Section~\ref{sec:cnl}).

\paragraph{Two audiences.}
L4 must satisfy two readers with incompatible habits. The machine wants an
unambiguous grammar with a denotational semantics. The lawyer wants prose,
or as close to prose as rigour permits, and has no patience for sigils. The
thesis of this paper is that these goals are reconcilable rather than
opposed: a language can have a wholly conventional formal semantics and,
\emph{at the same time}, a surface syntax tuned as a controlled natural
language. Sections~\ref{sec:constitutive}--\ref{sec:relational} make the
first case; Section~\ref{sec:cnl} makes the second.

% ===================================================================
\section{The Constitutive Core: A Typed Functional Language}\label{sec:constitutive}
% ===================================================================
L4's constitutive core is a statically typed, total, higher-order
functional language. Types are introduced with \code{DECLARE}: enumerations
and algebraic data types with \code{IS ONE OF}, and records with
\code{HAS}.

\begin{lstlisting}
DECLARE RiskCategory IS ONE OF
    LowRisk
    MediumRisk
    HighRisk
    Uninsurable

DECLARE Driver HAS
    name           IS A STRING
    age            IS A NUMBER
    accidentCount  IS A NUMBER
    hasTickets     IS A BOOLEAN
\end{lstlisting}

Functions are introduced by a signature---\code{GIVEN} naming the inputs and
\code{GIVETH} (synonym \code{GIVES}) naming the result---followed by a
definition with \code{MEANS}, or, for boolean-valued rules, by
\code{DECIDE\,\ldots\,IF}. Pattern matching uses \code{CONSIDER}\,/\,\code{WHEN},
a type-safe, exhaustiveness-checked case analysis.

\begin{lstlisting}
GIVEN driver IS A Driver
GIVETH A RiskCategory
`assess risk` driver MEANS
    CONSIDER driver's accidentCount
    WHEN 0 THEN IF driver's hasTickets
                THEN MediumRisk
                ELSE LowRisk
    WHEN 1 THEN MediumRisk
    WHEN 2 THEN HighRisk
    OTHERWISE   Uninsurable
\end{lstlisting}

Three properties of this core matter for law. \emph{Totality and
exhaustiveness}: the type checker rejects a \code{CONSIDER} that fails to
cover its cases, so a formalised rule cannot silently omit a category of
applicant---the analogue of an unhandled edge case is a compile-time error,
not a wrong benefit cheque. \emph{Type safety}: a payout formula cannot
accidentally add a date to a currency. \emph{Explainability}: because the
core is a pure functional evaluator, every result carries a trace from the
top-level question down to the deciding sub-condition, with each step
annotated by the rule that justified it. This trace is what the system shows
an end user, and what it hands to a language model as ground truth to
verbalise, rather than asking the model to reason about the law itself.

\paragraph{An isomorphic example.}
Consider a fragment of the British Nationality Act, the canonical AI-and-Law
benchmark \cite{sergot1986bna}. The statute's nested ``or''s and ``and''s
appear directly in the encoding, and the identifiers are the statutory
phrases themselves (backticks let an identifier contain spaces):

\begin{lstlisting}
GIVEN p IS A Person
GIVES A BOOLEAN
DECIDE `is a British citizen` IF
         `is born in the United Kingdom after commencement` p
      OR `is born in a qualifying territory after the appointed day` p
  AND       `is a British citizen`        OF `father of` p
         OR `is a British citizen`        OF `mother of` p
      OR    `is settled in the United Kingdom` OF `father of` p
         OR `is settled in the United Kingdom` OF `mother of` p
\end{lstlisting}

Note that the indentation, not parenthesisation, carries the precedence of
the boolean tree; we return to this in Section~\ref{sec:cnl}. Functions
marked \code{@export} are lifted automatically into a REST endpoint, an
OpenAPI schema, and a tool callable by an LLM agent; a single source file
thus yields an interactive eligibility wizard with citations back to the
rule that decided each branch.

% ===================================================================
\section{The Regulative Layer: Contracts as Concurrent Processes}\label{sec:regulative}
% ===================================================================
Definitional functions cannot express an obligation. ``The buyer must pay''
is not true or false of a situation; it is a demand on future conduct that
is met or breached as events unfold, possibly triggering a remedy. For this
L4 provides a regulative sub-language built from a five-keyword skeleton:

\begin{lstlisting}
PARTY   actor
MUST    action            -- or MAY, SHANT, DO
WITHIN  deadline
HENCE   successContinuation
LEST    failureContinuation
\end{lstlisting}

\code{MUST}, \code{MAY}, and \code{SHANT} are the deontic modalities
(obligation, permission, prohibition); \code{DO} is an uncommitted
action. \code{WITHIN} attaches a deadline. \code{HENCE} and \code{LEST} are
the continuations taken when the obligation is, respectively, satisfied or
not---and crucially, \emph{what counts as satisfaction depends on the
modality}: for \code{MUST}, performing the action by the deadline takes the
\code{HENCE} branch and letting the deadline lapse takes \code{LEST}; for
\code{SHANT} the polarity flips, so the prohibition is honoured precisely
when the deadline passes with no such action. \code{LEST} typically chains a
\emph{reparation}: a further obligation that repairs the breach, the
formal home of contrary-to-duty structure that bedevils Standard Deontic
Logic \cite{chisholm1963ctd,vonwright1951deontic}.

\begin{figure*}[t]
\begin{lstlisting}
DECLARE Person IS ONE OF Seller, Buyer
DECLARE Action IS ONE OF
    delivery
    payment HAS amount IS A NUMBER

saleContract MEANS
    PARTY  Seller
    MUST   delivery
    WITHIN 3
    HENCE ( PARTY  Buyer
            MUST   payment 100
            WITHIN 7
            HENCE  FULFILLED
            LEST   BREACH BY Buyer BECAUSE "payment not made within 7 days" )
    LEST   BREACH BY Seller BECAUSE "delivery deadline exceeded"

#TRACE saleContract AT 0 WITH
    PARTY Seller DOES delivery     AT 2
    PARTY Buyer  DOES payment 100  AT 5
-- Result: FULFILLED
\end{lstlisting}
\caption{A two-party sale as a regulative contract. \code{HENCE} threads the
seller's delivery into the buyer's payment obligation; \code{\#TRACE}
replays a timed event sequence and reports \code{FULFILLED}, a blame-bearing
\code{BREACH}, or the residual obligation that remains.}
\label{fig:sale}
\end{figure*}

\paragraph{Trace-based semantics.}
The denotation of a regulative contract is a function from a timed event
trace to an outcome: \code{FULFILLED}, a \code{BREACH} carrying blame and a
reason, or a \emph{residual} contract describing what remains owed. This is
exactly the trace-based model that Hvitved gives for multiparty contracts
in his Contract Specification Language \cite{hvitved2011thesis,
hvitved2012trace}, itself a descendant of the compositional commercial- and
financial-contract DSLs of Andersen, Elsman and Henglein and of Peyton
Jones and Eber \cite{andersen2006compositional, pjeber2000composing}. A
contract \emph{is} the set of event traces that conform to it; conformance
checking is trace membership; and the residual is the contract's state after
a prefix of events has been consumed---the machine-readable answer to ``what
is still owed, by whom, and by when.'' L4's \code{\#TRACE}
(Figure~\ref{fig:sale}) is precisely this fold of observed events through the
denotation. We borrow a great deal from CSL directly: its catalogue of
obligations, permissions, and prohibitions; its deadlines and its
reparation (contrary-to-duty) structure; and above all its foundational
stance that a contract simply \emph{is} the set of traces that conform to
it. To CSL's contract DSL L4 adds the two things it did not set out to
provide---a controlled-natural-language surface and an integrated
constitutive functional core---so that the predicates a contract branches
on are themselves type-checked, executable law rather than opaque
atoms. The dual view of a contract as a narrative of timed events that
initiate and terminate obligations is also exactly the setting of the event
calculus \cite{kowalski1986events}, which we take as the logical reading of
the same traces (Section~\ref{sec:relational}).

\paragraph{The CSP lineage.}
Read operationally, the regulative layer is a small process calculus in the
tradition of Hoare's Communicating Sequential Processes
\cite{hoare1978csp,hoare1985csp,roscoe1997theory}. The correspondence is
direct (Table~\ref{tab:csp}): an obligation \code{PARTY p MUST a HENCE k} is
a guarded event prefix $p.a \rightarrow k$; \code{HENCE} and \code{LEST}
form a choice resolved by what the environment---the trace---actually does,
in the manner of CSP's external choice; \code{RAND} is parallel composition
that synchronises on completion, so a compound breaches if either side does;
\code{ROR} is choice between sub-contracts; \code{FULFILLED} is successful
termination ($\mathit{SKIP}$) and \code{BREACH} is blame-assigning
termination. The one essential addition is \emph{time}: \code{WITHIN}
imposes a deadline, placing L4 with Timed CSP \cite{reed1988timedcsp} and the
timed-automata model that underlies tools such as UPPAAL
\cite{alur1994timed,behrmann2004uppaal}. Recursion gives recurring
obligations---a loan's monthly instalments are a function that re-emits next
period's obligation in its \code{HENCE} branch and a penalised obligation in
its \code{LEST} branch, bottoming out at \code{FULFILLED}.

\begin{table}[t]
\caption{L4 regulative keywords and their CSL / CSP analogues.}
\label{tab:csp}
\small
\begin{tabular}{@{}ll@{}}
\toprule
\textbf{L4} & \textbf{CSL / CSP analogue} \\
\midrule
\code{PARTY p MUST a HENCE k} & guarded prefix $p.a \rightarrow k$ \\
\code{WITHIN t}               & timed deadline (Timed CSP clock guard) \\
\code{HENCE k}                & success continuation \\
\code{LEST k'}                & failure / timeout cont.; reparation (CTD) \\
\code{A RAND B}               & synchronising parallel $A \parallel B$ \\
\code{A ROR B}                & choice $A \,\Box\, B$ \\
\code{FULFILLED}              & successful termination ($\mathit{SKIP}$) \\
\code{BREACH}                 & blame-assigning termination \\
\code{\#TRACE c WITH es}      & trace membership $es \in \mathrm{traces}(c)$ \\
\bottomrule
\end{tabular}
\end{table}

\paragraph{Why this matters.}
Because the regulative layer is a timed concurrent system, the questions a
lawyer asks about a contract become questions a verification tool can
answer. ``Is there a state in which one clause obliges an act that another
forbids?'' is a reachability query over the product of the parties'
processes---a race condition or, in deontic terms, a contradictory
conflict. In one of our public-sector deployments
(Section~\ref{sec:experience}), exactly such a deontic-temporal double bind
was discovered automatically in live regulation: under an unusual
interleaving, a party was simultaneously obliged and prohibited from the
same act. Recasting loophole-finding as state-space search is the bridge to
the relational reading we turn to next.

% ===================================================================
\section{Two Readings of One Program: Functional and Relational}\label{sec:relational}
% ===================================================================
A function $f : A \to B$ is, set-theoretically, its graph: the relation
$\{(a, f\,a) : a \in A\}$. Functional-logic languages exploit this
identity. Mercury annotates predicates with \emph{modes}---which arguments
are inputs and which outputs---and \emph{determinism}, then compiles the
relation to efficient code in whichever mode is requested
\cite{somogyi1996mercury}; Curry integrates functions and relations through
narrowing and residuation, so that a function may be run with unknown inputs
and the system searches for those that yield a desired output
\cite{antoy2010flp,hanus2013curry}. L4 inherits this perspective. Its
constitutive rules, being total typed functions, compile not only to a
forward evaluator but also to a relational form: Horn clauses in the manner
of Prolog, or, for the regulative layer, a transition relation over events
and fluents in the manner of the event calculus
\cite{kowalski1986events}. Logical English \cite{kowalski2021logicalenglish}
demonstrates that this declarative substrate can itself wear a
natural-language surface; L4 reaches it from the other direction, starting
from typed functional programming and recovering the relational reading by
transformation, with Mercury and Curry as the conceptual bridges between the
two paradigms.

This second reading is not a curiosity; it is what makes L4 a tool for
\emph{analysis} rather than mere execution. Three capabilities follow.

\emph{Backward / abductive queries.} Run forward, \code{`is a British
citizen`} classifies a given person. Run backward, the same rule
characterises the situations under which citizenship holds---useful for
explanation (``which facts, if changed, would flip this decision?'') and,
in the regulative layer, for planning: ``what is the \emph{latest} I may
apply for the loan and still deliver on time?'' is an abductive query over
the contract's transition relation, of the kind logic programming answers
natively but a one-directional evaluator cannot.

\emph{Exhaustive scenario generation.} Treating a rule as a relation lets a
property-based tester enumerate the space of situations and check
invariants across all of them, rather than the handful of hand-written
examples that legal code typically ships with. An early experiment imported
a national benefits calculation---originally a Python rules library with a
mere handful of unit tests---into a functional setting and generated tests
across the whole domain of scenarios, surfacing miscalculations that the
sparse hand-written suite had missed.

\emph{Model checking.} The transition relation can be handed to symbolic
model checkers---NuSMV, SPIN, TLA+, UPPAAL
\cite{cimatti2002nusmv,holzmann1997spin,lamport2002specifying,behrmann2004uppaal}---to
search for states that violate a property written at the specification
level. Here the letter of the law lives at the object level and the spirit
of the law lives as property assertions; a counterexample trace is a
loophole, found the way a security researcher finds an exploit. The
deontic-temporal conflict of Section~\ref{sec:regulative} is one such
counterexample.

L4's implementation already carries the seeds of this multi-target
strategy: a query planner reorders relational goals, and an MLIR/WASM
backend compiles decision logic to a portable binary
(Section~\ref{sec:impl}). The point of principle is that \emph{one}
isomorphic source, authored once by a domain expert, feeds all of these
analyses; the functional and relational readings are two views of the same
text, not two artefacts to be kept in sync.

% ===================================================================
\section{Surface Syntax as Controlled Natural Language}\label{sec:cnl}
% ===================================================================
Everything above concerns semantics, and any number of concrete syntaxes
could carry it. The wager of L4 is that the concrete syntax is where
adoption is won or lost, and that a formal language can be shaped into a
\emph{controlled natural language}---a restricted but rigorous fragment of
English \cite{kuhn2014cnl}---without sacrificing its formal character. The
design follows a few principles, each with a specific motivation.

\paragraph{Words, not sigils.}
L4 prefers familiar words to punctuation and symbols: \code{AT LEAST} for
$\geq$, \code{PLUS} for $+$, \code{'s} for field access (\code{driver's
age}), \code{OF} for application. A type signature reads as an English
clause---``\code{GIVEN} a person, \code{GIVES} a boolean''---rather than as
\code{Person -> Bool}. The cost is verbosity; the benefit is that a reader
who has never seen a type signature can still parse the sentence. This
directly serves the dual-legibility requirement of
Section~\ref{sec:motivation}. The genitive clitic \code{'s} is a small but
pleasing overload in this spirit: English writes \code{'s} both for
possession and as the contraction of \emph{is}, so \code{driver's age} reads
naturally as a \emph{has-a} (the field the driver \emph{has}) while the very
same clitic glosses as the copula in \emph{is-a} predications---one mark
spanning exactly the two relations that object orientation keeps apart as
composition and subtyping.

\paragraph{Two registers.}
The words that replace the sigils are set in upper case, and the choice is not
merely cosmetic: it splits the surface into two visual registers. \code{UPPERCASE}
tokens are the fixed vocabulary of the \emph{language}---\code{AND}, \code{OR},
\code{NOT}, \code{AT LEAST}, \code{GIVEN}, \code{MEANS}---while
\bq backtick-quoted lowercase phrases\bq{} are the vocabulary of the \emph{law},
the terms lifted from the source instrument. A reader can therefore tell at a
glance which words carry logical force and which merely name domain concepts,
without yet knowing the grammar---the same typographic convention by which COBOL
and SQL set their keywords apart from the programmer's own identifiers. The case
distinction does double duty as documentation.

\paragraph{Indentation, not parentheses.}
Following the offside rule \cite{landin1966next700}, L4 is
layout-sensitive, and it puts the layout to specific legal use: the
indentation of conjuncts and disjuncts encodes the precedence of the
boolean tree, so that the visual shape of the code is the logical shape of
the rule. In the British Nationality fragment of
Section~\ref{sec:constitutive}, the alignment of the \code{OR}s under and
beside the \code{AND} is what disambiguates the structure---no reader need
balance a parenthesis, and the indentation can be made to mirror the
sub-paragraph numbering of the source, advancing isomorphism. This is less the
import of a programming habit into law than the reverse: indented, enumerated
sub-paragraphs are how the written law has been set out since the Victorian
reforms of legislative drafting \cite{coode1843}, and an ordinary word processor
already indents a statutory list the same way---L4 only makes that indentation
carry the precedence.

\paragraph{Inert scaffolding and asyndeton.}
Two further devices let a formalisation reproduce its source almost word for
word. Where a statute joins its limbs without repeating a connective on every
line---the \emph{asyndetic} enumeration typical of drafted lists---L4 offers
\code{...} (three dots) and \code{..} (two dots) as lightweight infix synonyms
for \code{AND} and \code{OR}. The dot counts are deliberate: three for the three
letters of \code{AND}, two for the two of \code{OR}, so that in the fixed-width
register the implicit mark and its spelled-out form occupy the same column and
the limbs align. The connective can thus sit unobtrusively at the head of a
continuation line. And a quoted string in a boolean position is treated as
\emph{inert}: it evaluates to \code{TRUE} inside an \code{AND} and to
\code{FALSE} inside an \code{OR}, either way leaving the surrounding condition's
truth value untouched. Inert text lets the drafter paste the chapeau and
connecting prose of the source---``shall subsist for the life of the author
and'', ``whichever expires first''---directly into the rule as material the type
checker accepts and the evaluator ignores. Together with the layout, these let an
enumerated provision be transcribed limb for limb (Figure~\ref{fig:asyndeton}),
which is the surface condition for a lawyer to agree that the code is the law.

\begin{figure}[t]
\begin{lstlisting}
`grant is allowed` MEANS
        "the grant is allowed where all of the following hold---"
    ... `the form was filed on time`
    ... `the fee was paid`
    AND `the applicant is eligible`

`is eligible` MEANS
       "the applicant is eligible if any of---"
    .. `a citizen`
    .. `a permanent resident`
    OR `the holder of a valid pass`
\end{lstlisting}
\caption{Asyndeton and inert scaffolding. \code{...} is an implicit \code{AND}
and \code{..} an implicit \code{OR}, with the final limb spelling out the
conjunction as a drafted list does; the quoted chapeau is inert---it contributes
nothing to the truth value---and carries verbatim source text into the rule.}
\label{fig:asyndeton}
\end{figure}

\paragraph{Ditto marks.}
Tabular legal and financial text repeats structure down columns; clerks have
long abbreviated the repetition with the ditto mark. L4 borrows the
convention with a caret, \code{\textasciicircum}, which stands for the token directly above
it, so that aligned tabular data---a schedule of territories, an inventory,
a fee table---can be written the way a person fills in a form, with repeated
material elided into columns (Figure~\ref{fig:ditto}). The point is not
keystrokes saved but cognitive load: the eye reads the column as a unit.

\begin{figure}[t]
\begin{lstlisting}
myDeal MEANS
  Deal WITH
    buyer  IS Party WITH name IS "Alice Avocado"
    seller IS Party WITH name IS "Robert Rich"
    item   IS LIST InventoryItem OF "potato",  12,  Nothing
                                ^ OF "corn",    120, Nothing
                                ^ OF "avocado", 2,   Nothing
    price  IS Money OF usd, MoneyAmount OF 100, 00
\end{lstlisting}
\caption{The ditto caret \code{\textasciicircum} repeats the token above it, letting
aligned tabular data read as columns.}
\label{fig:ditto}
\end{figure}

\paragraph{Phrase-shaped identifiers and mixfix.}
Backtick-delimited identifiers may contain spaces---borrowing the convention,
familiar to every writer of Markdown, by which backticks fence a span of literal
text---so an identifier can be the exact statutory phrase
(\code{\bq is settled in the United Kingdom\bq}).
Such names may be \emph{mixfix}: arguments interleave with the name, so a
predicate reads as a sentence---\code{\bq employee\bq{} \bq works for\bq{} \bq employer\bq}
rather than \code{worksFor(employee, employer)}. Together these let the
formalisation quote its source, the surface condition for a lawyer to accept
that the code says what the law says (Figure~\ref{fig:mixfix}).

\begin{figure}[t]
\begin{lstlisting}
GIVEN `the worker` IS A Person
      `the firm`   IS A STRING
GIVES A BOOLEAN
`the worker` `works for` `the firm` MEANS
    `the worker`'s employer EQUALS `the firm`

DECIDE `Alice is entitled to leave` IF
    `Alice` `works for` "Acme"
\end{lstlisting}
\caption{Mixfix: the keyword \code{\bq works for\bq} interleaves with its two
arguments, so both the definition and the call site
(\code{\bq Alice\bq{} \bq works for\bq{} \bq Acme\bq}) read as English clauses
rather than as \code{worksFor(...)}.}
\label{fig:mixfix}
\end{figure}

\paragraph{Lineage.}
L4's surface design inherits from a substantial tradition of controlled and
computational language for rules. Its low-code ancestors are COBOL and SQL
(Section~\ref{sec:motivation}), which first wagered that an uppercase,
English-shaped surface could let non-programmers read business rules; the
controlled-language tradition since has refined that wager. From Grammatical
Framework
\cite{ranta2004gf} we take the discipline of treating a surface syntax as a
type-theoretical grammar---an abstract syntax with concrete realisations---
which points toward multilingual rendering of the same rule. From the
wide-coverage parsing tradition of C\&C and Boxer \cite{curran2007cc,
bos2008boxer} we take the goal, run in the opposite direction, of mapping
language to logical form, the ingestion problem an LLM now eases. From the
enterprise standards SBVR \cite{omg_sbvr} and DMN \cite{omg_dmn} we take the
demonstrated appetite among business analysts---not only programmers---for
structured vocabulary, rules, and decision tables; and from CNL efforts such
as RuleCNL \cite{feuto2014rulecnl}, which compiles a business-rule CNL into
SBVR, the specific lesson that a controlled surface can target a formal core
without exposing it. L4's distinctive bet relative to these is to pair the
CNL surface with an \emph{executable, type-checked} semantics spanning both
constitutive functions and regulative contracts, and to make the surface
\emph{isomorphic} to legal source text rather than to a business-analyst
vocabulary.

We do not claim these choices are free. Verbosity, layout-sensitivity, and
phrase identifiers each cost something in concision or in editor support,
and we discuss the trade-offs in Section~\ref{sec:experience}. We do claim
that they are the right trade-offs for a language whose first reader is a
domain expert, and that they are made \emph{at the surface}, leaving the
semantics of Sections~\ref{sec:constitutive}--\ref{sec:relational}
untouched.

% ===================================================================
\section{Implementation and Tooling}\label{sec:impl}
% ===================================================================
L4 is implemented in Haskell. The core (\textsf{jl4-core}) comprises a
layout-sensitive parser, a Hindley--Milner-style type checker, and an
evaluator that produces the explanatory traces of
Section~\ref{sec:constitutive}. Around it sits a tooling stack intended to
make a single \code{.l4} file productive: a Language Server Protocol
implementation for in-editor type checking, inline evaluation, and
go-to-definition; a REPL with trace and query-planning commands; a decision
service exposing \code{@export}-annotated functions as REST endpoints and as
Model Context Protocol tools for LLM agents; a WebAssembly build for
in-browser and serverless evaluation; and an MLIR backend that compiles
decision logic to a portable \code{.wasm} binary. The toolchain renders
every evaluation as a GraphViz \emph{ladder diagram}, and can transpile a
rule set to Markdown for human review or generate a consumer-facing web
wizard from the same source. A static analyzer classifies financial
contracts against the ACTUS and FIBO standards. The web editor at
\textsf{jl4.legalese.com} runs the entire pipeline client-side.

The role of LLMs in this stack is deliberately bounded. Models assist with
\emph{ingestion}---drafting a first L4 formalisation from a PDF or a statute
(Section~\ref{sec:ingestion})---and with \emph{explanation}---verbalising a
completed, deterministic evaluation trace for a lay user. They are not in the
reasoning loop. The reasoner is the L4 evaluator, whose every step is
auditable; the model's natural-language output is guard-railed by, and
checkable against, that trace. This is how L4 proposes to counter
hallucination in legal AI: not by trusting the model to reason about the
law, but by giving it a rigorous artefact to talk about.

% ===================================================================
\section{Ingestion: The Knowledge-Acquisition Bottleneck Revisited}\label{sec:ingestion}
% ===================================================================
% TODO: confirm the proper name and a citation for the "Legal Data Hunter" project.
The expert systems of the 1980s foundered on the
\emph{knowledge-acquisition bottleneck}
\cite{feigenbaum1984,hayesroth1983}: encoding a domain demanded a knowledge
engineer in slow, costly dialogue with a domain expert, and the brittle
result rarely repaid the labour. Legal expert systems were no exception. For
all the elegance of the British Nationality Act as a logic program
\cite{sergot1986bna}, formalising law one statute at a time, through a scarce
pairing of programmer and lawyer, never approached the scale of the law
itself. L4 attacks the bottleneck from two directions.

\paragraph{By hand, but better.}
Even hand encoding is easier in a language built as a controlled natural
language with isomorphism as a first-class goal
(Sections~\ref{sec:background} and~\ref{sec:cnl}). Because the formal text
mirrors the structure and quotes the wording of its source, a domain expert
can audit it clause by clause without becoming a programmer. We have found
this bridge to bear real weight: non-technical business stakeholders are
able to \emph{read} an L4 ruleset and \emph{sign off} on it as a faithful
statement of the rules they intend. This is a consequential result, because
sign-off is precisely the act by which legal responsibility passes from the
engineer who wrote the code to the accountable party who owns the
rules---the bridge that a formalism opaque to its domain experts can never
build. The knowledge-engineer/expert dyad does not vanish, but the gap
across it narrows; and, as our experience suggests, software engineers
schooled in requirements analysis are themselves capable wordsmiths who
surface the ambiguities a business expert would gloss over. Strong types and
exhaustiveness checking catch missing cases at the moment of encoding rather
than in production. What hand encoding cannot do is scale: it stays linear
in scarce expert hours.

\paragraph{At scale, with the type system as oracle.}
Large language models change the economics of the first pass. A model can
propose a formalisation of a statute or a policy in seconds, shifting the
human role from authoring to reviewing. The hazard is the familiar one---a
fluent encoding that is subtly wrong---and here L4's static type system and
language server are the decisive guardrail. The type checker is a sound,
deterministic oracle that rejects ill-typed and non-exhaustive
formalisations outright, and the LSP reports each violation inline with a
precise location, so a model can be driven in a generate--check--repair loop
until the file type-checks, and then until its \code{\#ASSERT} golden tests
pass---extending the oracle from well-typedness to behaviour. This is the
paper's neuro-symbolic thesis applied to ingestion itself: the LLM proposes
and the type system disposes. Strong typing turns open-ended
natural-language-to-code generation into a constrained search against a
verifier, the discipline that makes type-directed program synthesis
tractable, and it lets the system scale ingestion without surrendering
correctness to a model's confidence.

\paragraph{Corpus at scale.}
A type-guarded ingestion loop still needs raw material: machine-readable
primary law. The complementary enabler is the decades-long campaign to
liberate it. Carl Malamud's Public.Resource.Org and Law.gov, and the
litigation establishing that the law and the standards incorporated into it
are not locked behind copyright \cite{publicresource}, have made vast
quantities of primary legal material publicly accessible; projects in that
lineage---such as the Legal Data Hunter effort---are now channelling
statutes, regulations, and codes into ingestible form at scale. With an open
corpus on one side and a type-guarded LLM loop on the other, the isomorphic
formalisation of law begins to look tractable at a scale the
knowledge-acquisition bottleneck once foreclosed. It does not thereby become
automatic: faithfulness to legislative or contractual \emph{intent}, as
opposed to structure, remains a matter of human judgement and sign-off
(Section~\ref{sec:experience}). But formalisation need no longer proceed one
statute at a time, gated by a scarce dyad.

% ===================================================================
\section{Applications}\label{sec:apps}
% ===================================================================
Because a single \code{.l4} source feeds many backends
(Section~\ref{sec:impl}), one formalisation supports a spectrum of
applications across the legal lifecycle---drafting, analysis, citizen
self-help, and operational decisioning---rather than a single product. We
sketch four archetypes; the pilots of Section~\ref{sec:experience}
instantiate them.

\paragraph{Citizen-facing expert systems.}
From the same source that the type checker validates, the toolchain
auto-generates an interactive web wizard: a member of the public enters the
facts of their situation and is shown their obligations or entitlements,
each answer carrying a citation to the deciding rule and the evaluation
trace as a plain-language explanation. This is the disruptive wedge of
Section~\ref{sec:motivation}---the overserved non-consumer who simply wants
to know where they stand under a regulation or a policy they signed without
reading. Building such an expert system is, in L4, a side effect of having
written the rules down at all, not a separate engineering project.

\paragraph{Formal verification and static analysis.}
The relational reading (Section~\ref{sec:relational}) lets a rule set or a
draft contract be handed to model-checking backends---NuSMV, SPIN, TLA+,
UPPAAL---for compile-time static analysis: discovery of race conditions and
deontic-temporal double binds, of clauses that can never fire or
obligations that can never be discharged, and of outright inconsistencies in
legislation under drafting or in a contract under negotiation. Loophole
hunting becomes exploit hunting. The same machinery turns each party's
``dozen scenarios I care about'' into a regression suite that runs on every
redraft, so that an edit from either side that breaks a previously agreed
property is caught immediately rather than litigated later.

\paragraph{Generative drafting back into natural language.}
The ingestion direction---natural language to formal rule
(Section~\ref{sec:ingestion})---is the obvious use of an LLM, but the
reverse is the more interesting one. Once a rule set
is formal and has been \emph{proven consistent}, generative AI can render it
back into fluent English: the statutory prose a Parliament can debate and
enact, or the clause a counterparty will sign. Here the verified formal
artefact is the source of truth and the English is a faithful, regenerable
rendering---the inversion the rules-as-code agenda has long sought
\cite{mohun2020rulesascode}---rather than the prose being authoritative and
the logic an after-the-fact gloss that may silently disagree with it. Because
the model is constrained to verbalise a fixed artefact, the prose cannot
drift from the logic, and (via grammar-engineering approaches such as
\cite{ranta2004gf}) the same rule can be rendered in multiple natural
languages from one source.

\paragraph{Audit-grade enterprise decisioning.}
At the opposite, high-volume end, the decision service plugs into an
enterprise business-rules engine for operational decisions---approving or
denying loan applications, adjudicating insurance claims, determining
benefit eligibility. The differentiator is not throughput but
accountability: because every decision carries a deterministic trace from
the top-level question down to the deciding condition, the outcome is
explainable and auditable \emph{by construction}. This directly answers the
transparency expectations that the GDPR's provisions on automated decisions
gesture at \cite{gdpr2016}---and whose precise extent is contested
\cite{wachter2017explanation}---and that the EU AI Act makes concrete for
high-risk systems \cite{aiact2024}. The applicant is entitled to a reason,
and L4 produces the reason as a first-class artefact: a citable derivation,
not the post-hoc rationalisation of a model that, in effect, had a gut
feeling about them.

These archetypes are not four products but four views of one source; the
platform ambition (Section~\ref{sec:conclusion}) is precisely that an
ecosystem can build all of them, and others not yet imagined, atop rules
authored once.

% ===================================================================
\section{Experience and Open Questions}\label{sec:experience}
% ===================================================================
% TODO: confirm what may be disclosed about each engagement; figures and
%       counterparties are presently described in deliberately general terms.
L4 has been exercised in pilots across the public and private sectors. With
a government agency we encoded a piece of secondary legislation for
regulatory compliance and generated a web wizard that lets members of the
public enter their circumstances and see their obligations, each answer
carrying citations to the source rules; the formal-verification pass
uncovered the deontic-temporal conflict described in
Section~\ref{sec:regulative}. In the private sector, formalising a major
insurer's policy revealed an ambiguity in a payout formula with material
financial exposure. A legislative drafting office has explored L4 for
future drafting as a rules-as-code exercise of the kind the OECD has urged
\cite{mohun2020rulesascode}. A payments fintech used L4 to ingest commercial
agreements dense with fee tables and serve them through an SQL-like API for
integration with operational systems.

These engagements also exposed the language's rough edges. Layout
sensitivity, prized for isomorphism, can frustrate authors pasting from word
processors that mangle whitespace. Phrase-shaped identifiers strain editor
tooling built for token-shaped names. The regulative layer's defaults for
omitted \code{HENCE} and \code{LEST} clauses, while convenient, can surprise
authors who have not internalised the modality-dependent polarity of
``success.'' Each is a live design question. More fundamentally, the gap
between an isomorphic formalisation and a \emph{faithful} one---between
matching the structure of the text and matching the intent behind it---is
not something the language can close on its own; it remains, as it should, a
matter for human legal judgement, now better supported by tools.

% ===================================================================
\section{Related Work}\label{sec:related}
% ===================================================================
L4 sits at the confluence of several lines of research; its contribution is
their combination rather than any single strand.

\paragraph{Contract DSLs and process calculi.}
The financial-contract pearl of Peyton Jones and Eber
\cite{pjeber2000composing} and the commercial-contract algebra of Andersen
et al.\ \cite{andersen2006compositional} established the compositional,
combinator style that Hvitved's CSL \cite{hvitved2011thesis,hvitved2012trace}
generalised to multiparty contracts with a trace semantics and blame
assignment; L4's regulative layer is a direct descendant, adding a
controlled-natural-language surface and integration with a constitutive
functional core. Stipula \cite{crafa2022stipula} likewise models legal
contracts as a process calculus with state and time, reinforcing the
contracts-as-processes view from an automata angle.

\paragraph{Logic programming of law.}
The British Nationality Act project \cite{sergot1986bna} is the foundational
demonstration of isomorphic legislative encoding; PROLEG
\cite{satoh2011proleg} and the broader argumentation tradition
\cite{prakken2015law,bench2012history} continue it. Logical English
\cite{kowalski2021logicalenglish} shares L4's controlled-natural-language
ambition while remaining within logic programming; L4 differs in starting
from typed functional programming and reaching the relational reading by
transformation (Section~\ref{sec:relational}) rather than as its native
substrate.

\paragraph{Deontic and defeasible reasoning.}
Defeasible deontic logic and the Regorous line of work
\cite{governatori2005regorous,governatori2016semantics} model normative
positions, permissions, and contrary-to-duty structure with great
care. L4's reparation-via-\code{LEST} is a deliberately operational, less
expressive treatment of the same contrary-to-duty phenomena
\cite{chisholm1963ctd}, chosen for executability and for its fit with the
trace semantics.

\paragraph{Languages for law and smart contracts.}
Catala \cite{merigoux2021catala} is the closest contemporary in spirit: a
language for legislation, grounded in default logic, with a strong
correctness story. L4 differs in its explicit constitutive/regulative split,
its concurrency-aware contract semantics, and its CNL surface. DAML
\cite{daml} brings a Haskell-derived functional language to ledger
contracts but targets execution on distributed ledgers rather than
isomorphic formalisation of existing law. The symbolic-drafting tradition
runs back to Allen \cite{allen1957symbolic} and, earlier still, to Coode's
nineteenth-century anatomy of legislative sentences \cite{coode1843}.

% ===================================================================
\section{Conclusion and Future Work}\label{sec:conclusion}
% ===================================================================
We have presented L4, a functional programming language for computational
law that separates constitutive rules, treated as total typed functions,
from regulative rules, treated as timed concurrent processes with a
trace-based denotation in the lineage of CSL and CSP. We have shown that a
single isomorphic source admits both a functional reading, for evaluation
and explanation, and a relational reading, for abductive planning and
model-checking-based loophole detection; and we have argued that the
surface syntax can be engineered as a controlled natural language---low on
punctuation, layout-structured, phrase-named---so that the same text is
legible to the lawyer and to the machine. The motivation is not technical
novelty for its own sake but the durable demand, sharpened in the age of
generative AI, for rules a person can check for themselves precisely when
the law begins to bite.

Future work runs along three axes: deepening the formal correspondence with
CSL and Timed CSP to a full proof of adequacy for the regulative semantics;
maturing the relational backends so that conflict detection and abductive
planning are push-button for non-experts; and a controlled study of
authorability, measuring whether L4's CNL design in fact lowers the barrier
for legally trained, non-programmer authors. The broader aim is a platform:
a language at the bottom of the stack, in the manner of PostScript, on which
an ecosystem of legal applications---wizards, checkers, planners,
integrations---can be built atop rules written once, in black and white,
and read by everyone.

\begin{acks}
% TODO: funding, the Centre for Computational Law, collaborators, pilot partners.
The authors thank colleagues at the Centre for Computational Law and the
legal and engineering collaborators on the pilots described herein.
\end{acks}

% ===================================================================
\appendix
% ===================================================================
\section{Extending L4 to PROLEG: A Burden-of-Proof Monad}\label{app:proleg}
% TODO: this appendix describes work in progress (a parallel effort is building
%       the two-way parser/transpiler); soften claims further before submission.

This appendix sketches an extension of L4 toward \emph{PROLEG}
\cite{satoh2011proleg}, Satoh's logic-programming rendering of the
``presupposed ultimate fact theory'' of the Japanese Civil Code, and the
theory of the bridge we are building between the two.\footnote{The two-way
PROLEG$\leftrightarrow$L4 parser and transpiler are under active development
in a parallel effort; what follows is a design-and-theory note, not a report
on a finished system. The front end already round-trips canonical PROLEG
(\code{parse . print = id}) on the fixtures discussed below, and the
burden-of-proof library of \S\ref{app:burden} type-checks and reproduces the
reference judgement; the Mode-B emitter is in progress.} PROLEG is an
attractive target for three reasons aligned with the body of this paper: it
is a large, real corpus of formalised law (some 2{,}500 rules, validated
against decades of bar examinations), so it stresses ingestion at scale
(Section~\ref{sec:ingestion}); it is purely \emph{constitutive} and so
exercises L4's functional core and its relational reading
(Section~\ref{sec:relational}); and its native treatment of \emph{burden of
proof} forces a distinction the body of the paper left implicit---a third
jural role, neither constitutive definition nor regulative obligation.

\subsection{Canonical PROLEG}
Canonical PROLEG is ordinary Prolog augmented with a rule arrow \code{<=} and
a defeasibility operator \code{exception(H,E)}: rule \code{H} is defeated
when its defeater \code{E} is provable, and exceptions nest. Layered on top
is a litigation vocabulary---\code{allege}, \code{provide\_evidence},
\code{admission}, \code{plausible}---that encodes who must establish each
ultimate fact and to what standard. Figure~\ref{fig:proleg} shows the core of
the canonical lease fixture: a landlord's cancellation claim, an
approval-of-sublease defence, and an abuse-of-confidence exception that
defeats the tenant's \emph{non}-abuse defence---an exception of an exception.
The fixture is drawn from genuine Japanese law. Article~612 of the Civil
Code provides that a lessee may not assign or sublease without the lessor's
consent, and that upon an unauthorised sublease the lessor may cancel the
contract; the Supreme Court's judgement of 27~January 1966 (Minsh\=u 20-1-136)
qualified this with the breach-of-trust, or \emph{abuse of confidence},
doctrine, under which cancellation is barred where the sublease does not in
fact betray the parties' relationship of trust. That qualification is the
exception of an exception the fixture encodes; the original statutory text is
reproduced in the companion Japanese translation of this paper.

\begin{figure}[t]
\begin{lstlisting}[language=Prolog,basicstyle=\ttfamily\footnotesize,keywordstyle=\color{l4kw}\bfseries,commentstyle=\color{l4cmt}\itshape]
cancellation_due_to_sublease <=
    agreement_of_lease_contract, handover_to_lessee,
    agreement_of_sublease_contract, handover_to_sublessee,
    using_leased_thing, manifestation_cancellation.
exception(cancellation_due_to_sublease, get_approval_of_sublease).
exception(cancellation_due_to_sublease, nonabuse_of_confidence).
nonabuse_of_confidence <= fact_of_nonabuse_of_confidence.
exception(nonabuse_of_confidence, abuse_of_confidence).
abuse_of_confidence    <= fact_of_abuse_of_confidence.
\end{lstlisting}
\caption{Canonical PROLEG: the lease-cancellation rulebase
\cite{satoh2011proleg}, with two defeaters and an exception-of-exception.}
\label{fig:proleg}
\end{figure}

\subsection{Two modes of translation}
We translate in two modes that are, importantly, not two compilers.

\emph{Mode B (decision-only).} Erase the burden layer and read each rule
purely constitutively: \code{H <= B} becomes \code{DECIDE H IF B}, multiple
rules for one head become a disjunction, and each \code{exception(H,E)}
becomes a conjoined \code{AND NOT E}. Figure~\ref{fig:modeb} is the Mode-B
image of Figure~\ref{fig:proleg}. The translation is the classical reading of
stratified negation-as-failure \cite{clark1978naf,gelfond1988stable}, and is
sound exactly when the signed dependency graph is stratified---a condition
the transpiler must check, not assume.

\begin{figure}[t]
\begin{lstlisting}
DECIDE `cancellation due to sublease` IF
        `agreement of lease contract`
    AND `handover to lessee`
    AND `agreement of sublease contract`
    AND `handover to sublessee`
    AND `using leased thing`
    AND `manifestation cancellation`
    AND NOT `get approval of sublease`    -- exception
    AND NOT `nonabuse of confidence`      -- exception

DECIDE `nonabuse of confidence` IF
        `fact of nonabuse of confidence`
    AND NOT `abuse of confidence`         -- exception-of-exception
\end{lstlisting}
\caption{Mode-B L4: defeasibility becomes \code{AND NOT}; the
exception-of-exception falls out structurally. The output is plain
constitutive L4, validated by \code{jl4-cli}.}
\label{fig:modeb}
\end{figure}

\emph{Mode A (judgement-faithful).} Retain the burden of proof and reify it
as an L4 library. The key observation is that Mode~A is Mode~B \emph{plus}
typing each fact as a burden-bearing value and adding a \code{resolve} pass;
\code{resolve} is exactly the Mode-A-to-Mode-B desugaring. One compiler, two
fidelities.

\subsection{A clause-by-clause walkthrough}\label{app:walkthrough}
A translation earns the word \emph{isomorphic} only if a reader can lay the
two texts side by side and check them clause by clause. We walk the lease
case of Figure~\ref{fig:proleg} from rulebase to factbase to judgement, to
make the correspondence explicit.

\paragraph{Rules.} The mapping is local and structural. Two PROLEG rules that
share a head---\code{contract\_end <= cancellation\_due\_to\_sublease} and
\code{contract\_end <= expiration\_\ldots}---become one \code{DECIDE} whose
body is their disjunction. A rule together with the \code{exception/2} clauses
that name its defeaters---which PROLEG writes as \emph{separate} top-level
clauses---gathers into a single \code{DECIDE} whose body conjoins the rule's
premises with one \code{AND NOT} per defeater. This per-head gathering is the
sole normalisation the translation performs, and it is what makes the L4 rule
readable as a self-contained statement of when its head holds. Thus the
six-premise \code{cancellation\_due\_to\_sublease} rule and its two exceptions
collapse into the single declaration of Figure~\ref{fig:modeb}, and the chain
\code{nonabuse <= fact\_of\_nonabuse} / \code{exception(nonabuse, abuse)} /
\code{abuse <= fact\_of\_abuse} becomes the two declarations that close it.
Rule for rule, the L4 file has the same shape as the PROLEG one.

\paragraph{Facts.} The factbase (Figure~\ref{fig:leasefacts}) is where the
burden vocabulary does its work. A fact counts as \emph{established} when its
proponent meets the standard---\code{plausible(F)} holds---or when the
opposing party concedes it by \code{admission}. Reading the lease factbase
through that rule: the six constitutive facts of the cancellation claim are
\code{admission}-ed by the defendant, and so are established; the two approval
facts are \code{allege}d and evidenced by the defendant but never made
\code{plausible}, and so are \emph{un}established; while
\code{fact\_of\_nonabuse} and \code{fact\_of\_abuse} are each \code{plausible},
and so established. In Mode~B these resolve to the booleans the \code{DECIDE}s
consume; in Mode~A they are the \code{established}/\code{unestablished}
injections of Figure~\ref{fig:burdenlease}, each carrying its bearer into the
ledger.

\begin{figure}[t]
\begin{lstlisting}[language=Prolog,basicstyle=\ttfamily\footnotesize,keywordstyle=\color{l4kw}\bfseries,commentstyle=\color{l4cmt}\itshape]
admission(agreement_of_lease_contract, defendant). % +5 constitutive
allege(approval_of_sublease, defendant).
provide_evidence(approval_of_sublease, defendant). % but NOT plausible
allege(fact_of_nonabuse_of_confidence, defendant).
plausible(fact_of_nonabuse_of_confidence).         % => established
allege(fact_of_abuse_of_confidence, plaintiff).
plausible(fact_of_abuse_of_confidence).            % => established
\end{lstlisting}
\caption{The lease factbase (excerpt). \code{admission} by the opponent and
\code{plausible} both establish a fact; an \code{allege}d fact never made
\code{plausible} stays unestablished---and its bearer loses on it.}
\label{fig:leasefacts}
\end{figure}

\paragraph{Judgement.} Evaluating \code{contract\_end} now descends the rules
deterministically:
\begin{itemize}
\item \code{contract\_end} reduces to \code{cancellation} \code{OR}
  \code{expiration}; the expiration branch has no established facts and fails,
  so everything rests on cancellation.
\item \code{cancellation} requires its six constitutive facts (all established
  by admission), \code{AND NOT get\_approval}, \code{AND NOT nonabuse}.
\item \code{get\_approval} needs both approval facts, which are unestablished,
  so it fails---and \code{NOT get\_approval} holds. The defendant's first
  defence is out.
\item \code{nonabuse} needs \code{fact\_of\_nonabuse} (established) \code{AND
  NOT abuse}. But \code{abuse} needs \code{fact\_of\_abuse}, which \emph{is}
  established, so \code{abuse} holds, \code{NOT abuse} fails, and
  \code{nonabuse} fails---so \code{NOT nonabuse} holds. This is the exception
  of an exception: the plaintiff's abuse of confidence defeats the defendant's
  non-abuse defence.
\item With both defeaters out, \code{cancellation} holds, \code{contract\_end}
  holds, and the landlord prevails.
\end{itemize}
This descent \emph{is} the L4 evaluation trace---the ladder diagram of
Section~\ref{sec:impl}---and exactly what the \code{\#ASSERT}s of
Figure~\ref{fig:burdenlease} check. Because every step cites a rule that
exists one-for-one in the source, the trace doubles as an audit of the PROLEG
program: the explanation a reviewer reads is a literal walk over the statute's
own structure. The burden ledger then adds the dimension Mode~B
discards---\code{fact\_of\_abuse} on the plaintiff, \code{fact\_of\_nonabuse}
on the defendant---recovering, for free, the question of who would have lost
had each fact gone unproven.

\subsection{First-order rules and type reconstruction}
The lease fixture is propositional, but most of PROLEG is first-order:
predicates carry positional, untyped arguments. Figure~\ref{fig:minorduress}
shows a fragment of a second fixture---a minor's rescission of a sale,
defeated by duress---in which the arguments are buyers, sellers, contracts,
and dates. Translating it to L4, whose functions are typed and whose fields
are named, requires \emph{type reconstruction}: inferring from usage that the
first argument of \code{rescission\_by\_minor\_buyer} is the same sort as the
\code{Buyer} flowing into \code{minor} and \code{manifestation}, then giving
that position a name and a type. This is Hindley--Milner-style inference run
across the predicate signatures, and it is where untyped PROLEG and typed L4
most visibly diverge---the temporal arguments, for instance, become L4
\code{Date} values with a real ordering (Section~\ref{sec:regulative}), not
opaque atoms.

\begin{figure}[t]
\begin{lstlisting}[language=Prolog,basicstyle=\ttfamily\footnotesize,keywordstyle=\color{l4kw}\bfseries,commentstyle=\color{l4cmt}\itshape]
rescission_by_minor_buyer(Buyer,Seller,CID,Tc,Tr) <=
    minor(Buyer),
    manifestation(rescission(CID),Buyer,Seller,Tr),
    before_the_day(Tc,Tr).
exception(manifestation(A,Mfr,Mfe,Ta),
          manifestation_by_duress(Thr,Mfr,Mfe,A,Ta,Td,Tr)).
\end{lstlisting}
\noindent\textit{\footnotesize L4 (Mode B, types reconstructed):}
\begin{lstlisting}
GIVEN buyer IS A Person, seller IS A Person,
      contract IS A ContractId,
      tContract IS A Date, tRescission IS A Date
GIVETH A BOOLEAN
DECIDE `rescission by minor buyer`
       buyer seller contract tContract tRescission IF
        `minor` buyer
    AND `manifestation` (rescission contract) buyer seller tRescission
    AND tContract `before` tRescission

DECIDE `manifestation` action mfr mfe tAction IF
        `manifestation fact` action mfr mfe tAction
    AND NOT `manifestation by duress` ...   -- duress defeats it
\end{lstlisting}
\caption{First-order PROLEG and its Mode-B image. Positional untyped
arguments are reconstructed into named, typed parameters; the
exception-of-exception (duress defeats the minor's rescission, so the sale
stands) survives the translation.}
\label{fig:minorduress}
\end{figure}

Reconstruction is not merely a convenience: it is a \emph{linter}. The
published slides for this very example contain typos---a defeater spelled
\code{minifestation\_by\_duress}, variables \code{Maniester} and
\code{Mnifestee}---that untyped Prolog accepts silently, since a misspelt
functor is simply a fresh predicate that never unifies (a defeater that can
never fire, and so a bug that quietly changes the judgement). L4's name
resolver and type checker reject exactly these as unbound or ill-typed. The
type-as-oracle ingestion loop of Section~\ref{sec:ingestion} therefore pays a
second dividend here: round-tripping PROLEG's corpus through L4 makes the
corpus audit itself, surfacing the silent drafting errors that a large body
of untyped rules inevitably accretes.

\subsection{The burden-of-proof monad}\label{app:burden}
The recurring move in all three traditions the paper touches---deontic
obligation, residual control, and burden of proof---is to index a normative
position by a responsible party, $(\textit{position},\textit{party})$. By
Hohfeld \cite{hohfeld1913} these are three \emph{distinct} positions: the
obligation-bearer who is blamed on breach (a duty, the agentive subject of
CSL and of L4's \code{PARTY \ldots MUST}, Section~\ref{sec:regulative}); the
residual-control holder who decides in the gap (a power); and the bearer of
the risk of non-persuasion who \emph{loses} if a fact is unproven (a
liability). The construct here formalises \emph{only} the third, and the
governing transpiler rule is that a PROLEG party is a burden-role and must
\emph{never} be synthesised into an L4 obligation.

A proposition's evidential state has two independent dimensions: its
\emph{truth} (established, refuted, or undecided---naturally \code{Maybe},
with \code{Nothing} the undecided case) and its \emph{burden} (who must
establish it). ``A value carrying its evidential subject'' is the coreader /
environment comonad $(\textit{Subject}, a)$ \cite{uustalu2008comonadic}; it
becomes a \emph{monad} precisely when the subject carries a monoid---that is,
once one decides how two burdens combine. There is no sensible
$\textit{plaintiff} \mathbin{\diamond} \textit{defendant}$, so we do not
merge subjects but \emph{accumulate} them; the free monoid over elementary
burdens is an \emph{obligation ledger}, and the comonad-with-monoidal-output
is then a \code{Writer} monad \cite{wadler1992essence}. Stacked over the
truth dimension in the order that keeps the ledger even on a failed
proof---so that blame survives failure---this is
\code{MaybeT (Writer [Obligation])} (Figure~\ref{fig:provable}). The ledger
emitted as the \code{Writer} output \emph{is} the responsibility map,
generated for free by composition; and because conjunction is accumulating
rather than short-circuiting, that map is structural (it reflects the rule,
not the run).

\begin{figure}[t]
\begin{lstlisting}[language=Haskell,basicstyle=\ttfamily\footnotesize,keywordstyle=\color{l4kw}\bfseries,commentstyle=\color{l4cmt}\itshape]
data Subject    = Plaintiff | Defendant
data Obligation = Obligation { onWhom :: Subject, what :: Text }

-- truth (Maybe) over a ledger of burdens (Writer); MaybeT is
-- outside, so Provable a  =  ([Obligation], Maybe a)
type Provable = MaybeT (Writer [Obligation])

prove      :: Subject -> Text -> Maybe a -> Provable a -- a borne fact
notProven  :: Provable a -> Provable ()  -- negation as failure
flipBurden :: Provable a -> Provable a   -- involution = opposite(P)
absent     :: Provable a -> Provable ()  -- defeater, opposing party
resolve    :: Provable a -> Bool         -- close the world
\end{lstlisting}
\caption{The burden monad \code{L4.Proleg.Burden}. \code{flipBurden} is
PROLEG's \code{opposite(P)} localised to the exception boundary;
\code{absent = notProven . flipBurden} is a defeater the opposing party must
establish.}
\label{fig:provable}
\end{figure}

\subsection{The construct in L4 itself}\label{app:burdenl4}
The construct is not only Haskell theory; it is realised in L4 itself.
Figure~\ref{fig:burdenl4} excerpts \code{burden.l4}. Because L4 has no
\code{Monad} typeclass, the monad lives in the evaluator, and the same
algebra surfaces at the source level as a record carrying the truth
dimension (\code{MAYBE BOOLEAN}, with \code{NOTHING} the undecided case) and
an accumulating obligation ledger, together with explicit combinators.
\code{notProven} is negation as failure; \code{flipBurden} relabels a
sub-derivation's obligations to the opposite party; \code{resolve} closes the
world, so \code{NOTHING} becomes \code{FALSE}---the bearer loses. The
genitive accessor \code{p's truth}---the very \code{'s} overload of
Section~\ref{sec:cnl}---reads the record field.

\begin{figure}[t]
\begin{lstlisting}
DECLARE Subject IS ONE OF Plaintiff, Defendant

DECLARE Obligation HAS
    onWhom IS A Subject
    what   IS A STRING

-- Haskell: Provable = MaybeT (Writer [Obligation]).  L4 has no Monad
-- typeclass, so the algebra surfaces as a record: a truth dimension
-- (NOTHING = undecided) plus an accumulating obligation ledger.
DECLARE Provable HAS
    truth       IS A MAYBE BOOLEAN
    obligations IS A LIST OF Obligation

GIVEN p IS A Subject, claim IS A STRING, mx IS A MAYBE BOOLEAN
GIVETH A Provable
prove p claim mx MEANS
    Provable WITH
        truth       IS mx
        obligations IS LIST (Obligation WITH onWhom IS p, what IS claim)

-- negation as failure: holds iff its argument is not established
GIVEN p IS A Provable
GIVETH A Provable
notProven p MEANS
    Provable WITH
        truth       IS (IF isTrue (p's truth) THEN NOTHING ELSE JUST TRUE)
        obligations IS p's obligations

-- relabel obligations to the opposite party (an involution)
GIVEN p IS A Provable
GIVETH A Provable
flipBurden p MEANS
    Provable WITH
        truth       IS p's truth
        obligations IS flipAll (p's obligations)

-- close the world: established => TRUE, else the burden default
resolve p MEANS isTrue (p's truth)
\end{lstlisting}
\caption{Excerpt of \code{burden.l4}: the burden construct written in L4. The
record \code{Provable} is the L4 image of \code{MaybeT (Writer
[Obligation])}; the combinators are the source-level monad operations.}
\label{fig:burdenl4}
\end{figure}

The point of writing it in L4 rather than on paper is that the worked
judgement becomes executable. Figure~\ref{fig:burdenlease} encodes the lease
case of Figure~\ref{fig:proleg} directly and states the expected outcome as
\code{\#ASSERT}s that \code{jl4-cli} checks: \code{contractEnd} and
\code{cancellation} hold, while the approval defence and the non-abuse
defence do not---the latter defeated by the plaintiff's abuse of confidence,
an exception of an exception. The closing \code{\#EVAL} prints the
responsibility ledger; \code{flipBurden} dualises it, re-attributing each
fact to the party who bears its converse.

\begin{figure}[t]
\begin{lstlisting}
-- defendant's first defence: alleged but NOT plausible => unestablished
getApproval MEANS
  conj (LIST (unestablished Defendant "approval_of_sublease"),
             (unestablished Defendant "approval_before_cancellation"))

-- plaintiff's exception-of-exception
abuse MEANS conj (LIST (established Plaintiff "fact_of_abuse_of_confidence"))

-- defendant's second defence, itself defeated by the plaintiff's abuse
nonabuse MEANS
  conj (LIST (established Defendant "fact_of_nonabuse_of_confidence"),
             (notProven abuse))

cancellation MEANS
  conj (LIST (established Plaintiff "agreement_of_lease_contract"),
             -- ... the other five constitutive facts ...
             (established Plaintiff "manifestation_cancellation"),
             (notProven getApproval),
             (notProven nonabuse))

contractEnd MEANS disj (LIST cancellation)

#ASSERT      resolve contractEnd
#ASSERT      resolve cancellation
#ASSERT NOT (resolve getApproval)
#ASSERT NOT (resolve nonabuse)

#EVAL contractEnd's obligations          -- the responsibility ledger
#EVAL (flipBurden abuse)'s obligations   -- burdens flipped to the opponent
\end{lstlisting}
\caption{The lease judgement of \cite{satoh2011proleg} as executable L4. The
\code{\#ASSERT}s encode the reported outcome; the \code{\#EVAL}s print the
responsibility ledger that the \code{Writer} dimension accumulates.}
\label{fig:burdenlease}
\end{figure}

\subsection{The presumption is the monoidal identity}
Promoting the structure from a semigroup to a monoid forces a choice a
semigroup never makes: the identity element---the value of a proposition
\emph{before any evidence is combined in}. Legally, this unit is the
\emph{presumption}, and fixing it is a constitutional act. The carrier
\code{Bool} admits two monoids, De~Morgan duals
(Table~\ref{tab:presumption}): under $(\vee, \textit{False})$ a charge is not
established until a sufficient ground flips it true---the presumption of
innocence; under $(\wedge, \textit{True})$ a claim stands until disproof.
Same elements, same operation, \emph{different unit---different morality};
\code{flipBurden} De~Morgan-dualises the one into the other, mirroring
$\neg\neg = \mathrm{id}$. A \emph{rebuttable} presumption is the identity
(evidence can move off it); a \emph{conclusive} presumption or a \emph{jus
cogens} norm is the absorbing element (no evidence moves it). A legal regime
is thus a pair $(\textit{identity}, \textit{threshold})$---who gets the
benefit of the doubt, and how much doubt is tolerated; criminal law fixes
$(\textit{innocent}, \textit{beyond-reasonable-doubt})$, encoding
Blackstone's ratio \cite{blackstone1769} as, in effect, a decision-theoretic
loss function over the asymmetric cost of error \cite{hayspier1997burdens}.
The monad unit law reads jurisprudentially: the presumption on its own moves
nothing; only evidence ($\mathbin{>\!\!>\!\!=}$) does.

\begin{table}[t]
\caption{Burden of proof as the per-proposition choice of monoidal identity.}
\label{tab:presumption}
\small
\begin{tabular}{@{}lll@{}}
\toprule
\textbf{Monoid on \texttt{Bool}} & \textbf{identity} & \textbf{legal reading} \\
\midrule
$(\vee, \textit{False})$  & \textit{False} & presumption of innocence \\
$(\wedge, \textit{True})$ & \textit{True}  & rebuttable presumption of the claim \\
\bottomrule
\end{tabular}
\end{table}

\subsection{Correspondence and the answer-set image}
The same non-monotonic content has three faces the bridge must keep aligned
(Table~\ref{tab:correspondence}). Relationalising L4 as in
Section~\ref{sec:relational}---functions to predicates by A-normal-form
flattening with an added output argument---sends \code{Provable} and its
negation-as-failure to answer-set default negation
\cite{gelfond1988stable,dung1995argumentation}; stable-model enumeration
becomes scenario search, and the deontic-temporal double binds of
Section~\ref{sec:regulative} reappear as unsatisfiable cores, the same
loophole-as-exploit framing now extended to cover the burden layer.
Encoding the lease factbase and evaluating \code{contract\_end} with these
combinators (Figure~\ref{fig:burdenlease}) reproduces the reported
judgement: the landlord prevails, the
approval-of-sublease defence fails for want of plausibility, the non-abuse
defence is itself defeated by abuse of confidence, and the ledger attributes
each fact to its bearer---establishing, incidentally, that \emph{who bears a
burden} (the plaintiff, for the constitutive facts) is independent of
\emph{who establishes it} (here, the defendant's admission).

\begin{table}[t]
\caption{Three faces of the same defeasible content.}
\label{tab:correspondence}
\small
\begin{tabular}{@{}llll@{}}
\toprule
 & \textbf{PROLEG} & \textbf{\texttt{Provable}} & \textbf{ASP} \\
\midrule
defeater       & \code{exception} & \code{notProven}  & \code{not e} \\
closed world   & implicit         & \code{resolve}    & default neg. \\
proof standard & \code{plausible} & the \code{Maybe}  & weak constr. \\
burden bearer  & \code{opposite}  & \code{flipBurden} & tie-break \\
resp.\ map     & ---              & \code{Writer}     & literals/party \\
\bottomrule
\end{tabular}
\end{table}

The upshot for the main thesis is that the constitutive/regulative split of
Section~\ref{sec:background} gains a third lane---evidentiality. The natural
user-facing artefact is a \emph{responsibility map} with three projections
(who is on the hook, who decides, who must prove), each a
$(\textit{position},\textit{party})$ pair drawn from a different Hohfeldian
register, and the value of having them as separate types is precisely that
the transpiler cannot silently collapse one into another.

\subsection{Discharging the burden over time: proof actions and an event-sourced factbase}\label{app:proofaction}
The construct so far is static: it resolves a \emph{fixed} factbase to a
judgement. But a burden is discharged \emph{over time}, by acts performed before
deadlines---which is the regulative layer's home ground
(\S\ref{sec:regulative}), not the constitutive core's. This suggests siting the
\emph{act of proving} as a deontic action while leaving the \emph{fact} itself
constitutive. The modality must be \code{MAY}, not \code{MUST}: proving is a
privilege, not a duty---a party is not sanctioned for failing to prove, it
simply loses the point---so the Hohfeldian separation above survives even as the
deontic machinery, which carries a party index, is reused to name the
burden-bearer. A \emph{proof action}
\begin{lstlisting}
PARTY  Plaintiff
MAY    `prove` `abuse of confidence`
WITHIN evidenceDeadline
HENCE  `abuse of confidence is established`
\end{lstlisting}
\noindent feeds the constitutive decision: its \code{HENCE} continuation records
the fact a downstream \code{DECIDE} consumes. This splits the burden into its two
classical halves---the deadline-bounded \code{MAY} is the burden of
\emph{production} (act in time or lose the point), and the standard of proof
gating the \code{HENCE} is the burden of \emph{persuasion}---and renders a
defeater as a \emph{counter} proof action by the opposing party, ordered by
deadline, so that the static \code{flipBurden} involution becomes literal
turn-taking: claim, defence, reply, the shape of pleading itself.

The question this raises is not syntactic---\code{HENCE} already carries an
arbitrary expression---but semantic: how does the \code{HENCE} record a fact a
later \code{DECIDE} will read? The answer is \emph{not} a mutable cell and an
\code{UPDATE} keyword. Establishing facts is state that must never be
overwritten, for two reasons that pull together. The first is \emph{epistemic
revision}: evidence is found forged, a witness recants, an appeal admits new
facts, so a fact once believed established may later be believed false. On top of
the \emph{logical} non-monotonicity already present---the \code{AND NOT}
defeaters, which flip a conclusion \emph{within} a knowledge state---this adds a
second, \emph{temporal} non-monotonicity \emph{across} knowledge states. The
second reason is audit: the wedge of Section~\ref{sec:apps}, a decision one can
replay and explain, demands answering \emph{what did the system conclude on date
$T$, and on what basis}---and a destructive update is precisely what makes that
question unanswerable.

Both point to an \emph{event-sourced}, \emph{bitemporal} factbase. Nothing is
mutated: a proof action appends an event stamped with both \emph{valid time}
(when the fact holds in the world) and \emph{transaction time} (when we recorded
believing it) \cite{snodgrass2000temporal}, and a correction is a later,
superseding event rather than an overwrite. The facts a decision reads are then a
pure projection, \code{established(F) AS OF (validTime, knownAt)}, non-monotone
in either axis yet computed without mutation---and L4's multi-axis
\code{TemporalContext} (\S\ref{sec:regulative}) already supplies exactly these
stamps. The burden ledger of Figure~\ref{fig:provable} is already append-only;
all that changes is that \code{resolve} graduates from a fold into a temporal,
supersession-aware projection.

This is \emph{event sourcing}, and for this paper's audience it has a familiar
logical name: the \emph{event calculus} \cite{kowalski1986events}, whose
\code{HoldsAt(f,t)} is exactly \code{established(f) AS OF t} and from which the
regulative layer already descends (\S\ref{sec:regulative}). The same object
wears several hats---event sourcing for the engineer, the event calculus for the
logician, a \emph{trace} for the process reading of \S\ref{sec:regulative}, and,
for the lawyer, \emph{the record}: append-only, corrections filed rather than
expunged, interrogable as of any date. The capability L4 would need is therefore
not mutation but its dual---letting a constitutive \code{DECIDE} \emph{query the
regulative trace}---and we leave it to the work in progress flagged at the head
of this appendix.

\bibliographystyle{ACM-Reference-Format}
\bibliography{l4-icail}

\end{document}
